\section{Introduction}\label{sec:introduction}

Community recovery in the stochastic block model (SBM) is a basic problem in network statistics. The model assigns each node to an unknown community and makes edge probabilities depend on the communities of their endpoints \cite{holland1983}. Two common approaches use different representations of the same data. Spectral methods reduce the adjacency matrix to a small number of eigenvectors and then cluster the embedded nodes. Likelihood methods compare how well alternative community assignments explain the observed connections. We show that likelihood-ratio information can be decoded from the retained spectrum to achieve nearly optimal community recovery for every fixed $K$ under uniform signal conditions and $p=\Omega(\log n/n)$, where $p=\max_{a,b}P_{ab}$, without returning to the original adjacency rows.

This question concerns the decoding rule as well as the embedding. A low-rank approximation is selected by a quadratic criterion, and standard spectral clustering applies a further Euclidean criterion through $K$-means. The oracle classifier for an SBM instead weights connections by logarithms of block probabilities and includes the contribution of absent edges. For a general connectivity matrix and unequal community sizes, its decision boundaries need not agree with nearest-center boundaries in the spectral embedding. Nevertheless, the fact that spectral reduction uses a quadratic objective does not itself imply that it discards the relevant likelihood information: a projection preserves every linear statistic whose weight vector belongs to the projected subspace.

\paragraph{Optimal recovery and likelihood refinement.}
The information-theoretic limits of the SBM are well developed. Abbe and Sandon \cite{abbe2015} characterize exact recovery in the general logarithmic-density SBM through Chernoff--Hellinger divergence and give efficient algorithms reaching that threshold. Zhou and Li \cite{zhouli2020} study the optimal error rate through Chernoff bounds and its application to community detection. We use these works to identify the oracle benchmark, while the statistical objective of the present paper is the leading-exponent form
\[
    \exp\{-(1+o(1))I\},
\]
where $I$ is the relevant pairwise Chernoff information. Throughout this paper, ``nearly optimal'' means attaining the leading oracle Chernoff exponent. We do not claim the refined multiplicative rate or a new global minimax lower bound.

Spectral initialization followed by likelihood refinement is also an established approach. Amini, Chen, Bickel and Levina \cite{amini2013} fit mixtures to blockwise connection counts using pseudo-likelihood and iterative updates. Gao, Ma, Zhang and Zhou \cite{gao2017} prove that a weakly consistent initializer followed by penalized local likelihood refinement achieves optimal misclassification performance under their assumptions. Zhou and Amini \cite{zhouamini2020} obtain optimal bipartite recovery using spectral initialization and pseudo-likelihood classification, and analyze a broader family of EM-type biclustering updates. A recent preprint of Wang \cite{wang2026vem} analyzes spectrally initialized variational EM in a general SBM, obtaining a leading Chernoff exponent while estimating parameters and reusing the original graph. In these methods the refinement uses information calculated from the original edge observations. Our algorithm imposes an additional restriction: once the selected adjacency eigenpairs have been computed, initialization, parameter updates, classification and component replacement use only those eigenpairs and quantities reconstructed from them.

\paragraph{Likelihood information in spectral representations.}
Spectral optimality is known in important special cases. Abbe, Fan, Wang and Zhong \cite{abbe2020} establish a first-order entrywise approximation for eigenvectors and show that, in the equal-sized symmetric two-community SBM at logarithmic density, an untrimmed spectral method attains the optimal recovery threshold and misclassification exponent. Their result already demonstrates that a spectral representation can retain the decisive local evidence. For homogeneous planted-partition SBMs with multiple, possibly unequal communities, Zhang \cite{zhang2024} gives matching error exponents for high-degree trimming followed by eigenvalue-weighted spectral clustering and $K$-means. His spectral exponent is distinguished from the information-theoretic exponent, so a sharp analysis of a spectral algorithm does not automatically establish its statistical optimality.

Probabilistic inference based on spectral embeddings predates the present work. Suwan et al.\ \cite{suwan2016} use Gaussian-mixture limits of adjacency spectral embedding to construct an empirical Bayes prior for block membership inference. This is a relevant example of using distributional information beyond Euclidean clustering; it does not establish a sparse large-deviation equivalence between our spectral scores and the oracle edge likelihood ratio. For lower densities, Zhou and Amini \cite{zhouamini2019} provide data-driven adjacency regularization and consistency analysis in the general bipartite setting. These consistency results do not by themselves prove preservation of the oracle error exponent.

\paragraph{Our approach.}
We retain the $K$ eigenpairs of the original adjacency matrix with largest absolute eigenvalues, including signal eigenvalues of either sign, and form clipped reconstructed connection profiles $q_i$. We compare profiles through the complete Bernoulli cross-entropy, including the nonedge term. A spectral estimate of the probability scale sets the logarithm floor. The resulting weighted-mean M-step has the same form at sparse and dense probabilities. Keeping the complete score is necessary for our density range: a dense two-community example in Section~\ref{sec:main-results} shows that the sparse score can lose a fixed fraction of the oracle exponent even when the parameters are known.

The growing initialization starts a component at the global profile mean, adds a node-profile candidate according to its likelihood gain over all nodes, and refits every component after each addition. Algorithm~1 then performs $\lceil\log n\rceil$ rounds of single-component replacement. In each round it tests every deletion position, chooses the best observed node candidate for that deletion by the same all-node gain, and fits the resulting profiles from uniform weights. The old fit and all trial endpoints compete under one working mixture likelihood. One final E-step and parameter M-step precede classification. This is an explicit repeated version of the existing replacement operation; it uses no randomized restart.

A short population argument is the key to the general-$K$ proof. Assign each of the $K$ true profiles to its nearest fitted profile. If a fitted profile is unused, it can be replaced without harming any assigned class. Otherwise the assignment is a bijection, and the fitted profile assigned to the class with largest loss can be replaced. In either case, inserting that true profile removes at least $1/K$ of the population hard loss. Every true community supplies a good observed candidate on the spectral approximation event. The exhaustive deletion scans therefore transfer this contraction to the empirical objective, up to vanishing relative error and the $n\log K$ cost of passing from hard to soft loss. Repetition makes the soft loss $o(n^2p)$, which certifies weak responsibilities simultaneously for all $K$ communities. This argument does not require the growing path to preserve a pure center at every intermediate stage or any EM fit to reach a global optimum.

The likelihood interpretation has a precise algebraic basis. If $\widehat\Pi$ is the retained empirical projector, the signed rank-$K$ reconstruction equals $A\widehat\Pi$. Under a population SBM projector, every block-constant Bernoulli contrast is preserved when $P$ has full rank. Empirical projection, clipping and parameter estimation introduce errors. We control them on a full-graph event whose failure probability is at most $C_H\exp(-Hnp)$ for every fixed $H$. This confidence level matters when the information scale exceeds $\log n$. First-order entrywise eigenspace analysis, a direct clipped Chernoff moment bound, and a deterministic invariant-basin argument handle the dependence of the fitted profiles on the same graph.

\paragraph{Contributions and scope.}
Theorem~\ref{thm:main-recovery} proves an end-to-end bound
\[
 \mathbb E\,\Mis(\widehat z,z)\le e^{-(1+o(1))I_n}
\]
for the complete deterministic algorithm and every fixed $K$. The assumptions are positive community proportions, comparable block probabilities, a uniform singular-value lower bound relative to $p$, probabilities bounded away from one, and $p=\Omega(\log n/n)$. Under comparability, this density condition implies expected degrees $\Omega(\log n)$. The assumptions allow heterogeneous degrees, intermediate densities and dense graphs, and neither $P/p$ nor the proportions need have limits. The theorem includes initialization and fitted parameters. It yields exact recovery when $\liminf I_n/\log n>1$ and concerns the leading oracle exponent, not a refined multiplicative rate.

The deterministic replacement theorem also identifies which algorithmic change closes the initialization gap. It applies to any starting fit in the observed profile box, including a poor growing endpoint. We do not claim a general-$K$ theorem for the unchanged growing-only algorithm. A certified five-community population example shows why a claim covering arbitrary finite intermediate EM budgets would fail: one update at each growing stage can leave a missing community and a duplicated center even at an arbitrarily large signal scale. One replacement round repairs that example; it does not establish failure of fully converged growing in general.

Finally, a reproducible 60-graph benchmark evaluates the complete Bernoulli implementation for $K=2,\ldots,6$ at logarithmic and dense probability scales, with the same reconstructed profiles used for sparse-score controls. Replacement does not change the classification error in this collection, and weak general signals remain difficult. The earlier paired 240-graph experiments are reported separately because they evaluate the historical sparse-score growing and repair procedures. These numerical results expose finite-sample behavior and implementation properties without substituting for the theorem.

Section~\ref{sec:algorithm} specifies the algorithm, Section~\ref{sec:main-results} states the recovery results, and Section~\ref{sec:experiments} reports experiments. Proofs are in the appendix.
